\section{Even orders: supports with an even element}
\label{sec:notallodd}

Throughout this section the order $k$ is even. With \autoref{thm:sigmabump},
the following theorem disposes of every support at even order, other than the
full support, that is not all-odd. Reads are described in the vocabulary of
\autoref{sec:nao-reading}, in particular \eqref{eq:read-reads} and
\eqref{eq:read-pins}, and \autoref{lem:nao-mirror} turns reads of $1$ into
reads of $0$.

\begin{theorem}
\label{thm:notallodd}
Let $k$ be even and let $S\subseteq\{1,\dots,k-1\}$ be primitive and
palindromic, not the full support, and containing an even element. Then
$S$ carries a $\sigma$-bump $(\sigma,P)$.
\end{theorem}

The natural carriers at even order are the \emph{parity carriers}: $\sigma$
together with every free residue that has the parity of $\sigma$. Because $k$ is even, the reflection $y\mapsto k-y$ preserves
parity, and the argument splits along the two parity classes. Under (B1) the
reads of the pins are fixed by \eqref{eq:read-pins}, and a parity carrier puts every free position of one class
into $Q$, so an unwrapped tap that lands on a free position of that class
reads $1$. A clause can therefore fail, for $\sigma$ or (by a reflection) for
$k-\sigma$, only at an index $i$ where every unwrapped tap of the right class
lands back in $S$, so that $S$ is almost closed under adding $i$. If
$a=\min S$ is even, a failure of the odd class forces a gap in the even part
of $S$, and taps reading into that gap make the even class work for every
even $\sigma$ (\autoref{thm:nao-evenmin}). If $a$ is odd and both classes
fail for every $\sigma$, the failures force a rigid progression structure
(\autoref{thm:nao-struct}), which carries one of two explicit carriers at
$\sigma=a$ (Lemmas~\ref{lem:nao-bfone} and~\ref{lem:nao-bfshadow}). On a
first reading one can take Lemmas~\ref{lem:nao-noband},
\ref{lem:nao-endend} and~\ref{lem:nao-endband} on trust.

\subsection{Parity carriers}
\label{sec:nao-parity}

Write $S_0$ and $S_1$ for the even and odd elements of $S$, and refer to the
residue of an integer modulo $2$ as its \emph{class}. Since $k$ is even,
$s\equiv k-s\pmod2$ for every integer $s$, so $S_0$ and $S_1$ are each
palindromic, and a nonempty class $S_c$ has largest element $k-\min S_c$.
If $S$ is primitive then $S_1\ne\emptyset$, since otherwise $2$ would divide
every element of $S\cup\{k\}$.

\begin{definition}
For $\sigma\in S$, the \emph{parity carrier} of $\sigma$ is
\[
Q_\sigma=\{\sigma\}\cup\{\,y\in\{1,\dots,k-1\}:\ y\notin S,\ y\equiv \sigma
\pmod2\,\}.
\]
An index $i$ is \emph{silent} for $Q_\sigma$ if no surviving tap reads $1$
at $i$. For an index $i$ and a class $c$, write $C_c(i)$ for the statement
\begin{equation}
j\in S,\quad i+j<k,\quad i+j\equiv c\pmod2\quad\Longrightarrow\quad i+j\in S .
\label{eq:nao-cont}
\end{equation}
\end{definition}

\begin{lemma}[reads of a parity carrier]
\label{lem:nao-parity}
Let $k$ be even, $S$ palindromic, $\sigma\in S$, and let $c$ be the class of
$\sigma$.
\begin{enumerate}
\item[(a)] $Q_\sigma\cap S=\{\sigma\}$, so (B1) holds.
\item[(b)] $y\in Q_{k-\sigma}\iff k-y\in Q_\sigma$ for $0<y<k$.
\item[(c)] A surviving tap whose position has residue $y$ reads $1$ if and
only if either it is unwrapped, $y$ is free and $y\equiv c$; or it is wrapped and
$y\in S$ or $y\not\equiv c$.
\item[(d)] If $(\sigma,Q_\sigma)$ is not a $\sigma$-bump, then for
$\sigma'=\sigma$ or $\sigma'=k-\sigma$ (so $\sigma'\in S$ and
$\sigma'\equiv\sigma$) some index is silent for $Q_{\sigma'}$.
\item[(e)] If $i$ is silent for $Q_\sigma$, then $C_c(i)$ holds, and every
wrapped tap at $i$ reads a position with residue in $Q_\sigma$.
\end{enumerate}
\end{lemma}

\begin{proof}
(a) is the definition. (b) holds because $y=k-\sigma\iff k-y=\sigma$, because
$y$ is free exactly when $k-y$ is (palindromy), and because $y\equiv k-y$
and $k-\sigma\equiv\sigma$.
(c) is \eqref{eq:read-reads}, using that a surviving tap has $y\ne\sigma$.

(d) By (a), the pair fails (B2) at some index $i$. If no surviving tap reads
$1$ at $i$, take $\sigma'=\sigma$. Otherwise no surviving tap reads $0$ at
$i$, and then by (b) and \autoref{lem:nao-mirror} no surviving tap reads $1$
at the index $k-i$ for $Q_{k-\sigma}$; take $\sigma'=k-\sigma$. In both
cases $\sigma'\equiv\sigma$ since $k$ is even.

(e) Let $j\in S$ with $i+j<k$ and $i+j\equiv c$. If $i+j=\sigma$ then
$i+j\in S$. Otherwise the tap $j$ survives, and since it does not read $1$,
(c) says that $i+j$ is not free. For a wrapped tap whose position has residue
$y$: if $y=\sigma$ then $y\in Q_\sigma$; otherwise the tap survives, and since it
does not read $1$, $y$ is free and $y\equiv c$.
\end{proof}

For the rest of the section, $S$ is palindromic at even $k$ with an element
of each class, $a=\min S$, and $b$ is the least element of the class not
containing $a$. Then $\max S=k-a$, the class of $b$ has largest element
$k-b$, and $b\le k-b$. Since also $a<b$,
\begin{equation}
2a<k,\qquad a+b<k .
\label{eq:nao-ab}
\end{equation}

The next three lemmas say where a silent index can be. Near the start, the
taps $k-a$ and $k-\min S_c$ read unwrapped positions above the largest
element of their class, which are free.

\begin{lemma}[the top taps]
\label{lem:nao-ends}
Let $\sigma\in S$, of class $c$. An index $i$ is not silent for $Q_\sigma$
if either
\begin{enumerate}
\item[(i)] $i<a$ and $i\equiv a+c\pmod2$; or
\item[(ii)] $i$ is even and $i<\min S_c$.
\end{enumerate}
\end{lemma}

\begin{proof}
(i) The tap $k-a$ reads the position $y=i+k-a$, with $k-a<y<k$. So it is unwrapped,
$y$ is free because it exceeds $\max S$, and $y\equiv i+a\equiv c$. Being
free, $y\ne\sigma$, so the tap survives, and it reads $1$ by
\autoref{lem:nao-parity}(c).
(ii) Put $m=\min S_c$. The tap $k-m$ reads the position $y=i+k-m$, with $k-m<y<k$, and
$y\equiv m\equiv c$ since $i$ and $k$ are even. As $y$ is of class $c$ and
exceeds $\max S_c=k-m$, it is free, and the tap reads $1$ as before.
\end{proof}

Beyond $a$ the wrapped taps take over. Every element of $Q_\sigma$ has the
class of $\sigma$, and the wrapped taps $k-a$ and $k-b$ then confine a silent
index to one parity and to the interval $(a,b]$.

\begin{lemma}[the parity dichotomy]
\label{lem:nao-dichotomy}
Let $\sigma\in S$, of class $c$, and let $i>a$ be silent for $Q_\sigma$.
Then $i\equiv a+c\pmod2$ and $i\le b$.
\end{lemma}

\begin{proof}
All elements of $Q_\sigma$ have class $c$. The tap $k-a$ wraps at $i$
(as $i>a$), and its position has residue $i-a$, so by
\autoref{lem:nao-parity}(e) $i-a\equiv c$, that is $i\equiv a+c$. If $i>b$,
the tap $k-b$ also wraps, and in the same way $i\equiv b+c$; but
$b\not\equiv a$.
\end{proof}

Finally, a containment $C_c(i)$ at an odd index empties an interval of the
other class.

\begin{lemma}[the gap]
\label{lem:nao-gap}
Let $i$ be odd with $C_c(i)$, and suppose $S_c\ne\emptyset$. Then $S$ has no
element of class $1-c$ in the open interval $(i,\ i+\min S_c)$.
\end{lemma}

\begin{proof}
Let $s\in S$ have class $1-c$ and $s>i$. Then $k-s\in S$ has class $1-c$,
$i+(k-s)<k$ and $i+(k-s)\equiv c$ since $i$ is odd. By $C_c(i)$,
$i+k-s\in S_c$, so $i+k-s\le\max S_c=k-\min S_c$, that is $s\ge i+\min S_c$.
\end{proof}

\subsection{Even minimum}

\begin{theorem}[parity carriers at even minimum]
\label{thm:nao-evenmin}
Let $k$ be even and $S$ primitive and palindromic with $a=\min S$ even. If
$(\sigma_1,Q_{\sigma_1})$ is not a $\sigma$-bump for some odd
$\sigma_1\in S$, then $(\sigma_0,Q_{\sigma_0})$ is a $\sigma$-bump for every
even $\sigma_0\in S$. In particular $S$ carries a $\sigma$-bump on a parity
carrier.
\end{theorem}

\begin{proof}
Since $S$ is primitive, $S_1\neq\emptyset$, so $b=\min S_1$ exists, and
$a<b$. By \autoref{lem:nao-parity}(d) some index $i_1$ is silent for
$Q_{\sigma'}$, with $\sigma'\in S$ odd. \autoref{lem:nao-ends} with $c=1$
excludes $i_1<a$ odd (by (i), as $a+1$ is odd), and $i_1\le a$ even (by
(ii), as $\min S_1=b>a$). So $i_1>a$, and \autoref{lem:nao-dichotomy} gives
$i_1$ odd and $i_1\le b$. By \autoref{lem:nao-parity}(e) $C_1(i_1)$ holds,
and \autoref{lem:nao-gap} gives
\begin{equation}
S_0\cap(i_1,\ i_1+b)=\emptyset .
\label{eq:nao-evengap}
\end{equation}

Now fix an even $\sigma\in S$. I claim that every index $i$ has a surviving
tap reading $1$ for $Q_\sigma$. There are four cases.

\emph{$i<a$, $i$ even.} \autoref{lem:nao-ends}(i) applies, as $i\equiv a$.

\emph{$i<a$, $i$ odd.} The tap $b$ reads the unwrapped position $i+b<a+b<k$. The
position $i+b$ is even and $i_1\le b<i+b<a+b<i_1+b$, so by
\eqref{eq:nao-evengap} it is not in $S$. It is therefore free, of class
$0$, and not $\sigma$, and the tap reads $1$.

\emph{$i=a$.} Suppose $a$ is silent. Then $C_0(a)$ holds by
\autoref{lem:nao-parity}(e). Let $s$ be the largest element of
$Z=\{y\in S_0:\ y+i_1<k\}$, which contains $a$ because $a+i_1\le a+b<k$. By
$C_1(i_1)$, $s+i_1\in S$, and it is odd, so $s+i_1\le\max S_1=k-b$. Then
$a+s\le k-b-(i_1-a)<k$, so $C_0(a)$ puts $a+s\in S_0$; and
$a+s+i_1\le a+k-b<k$. So $a+s\in Z$, against the maximality of $s$.

\emph{$i>a$.} Suppose $i$ is silent. By \autoref{lem:nao-dichotomy} $i$ is
even and $i\le b$, so $i<b$ since $b$ is odd; and $C_0(i)$ holds. Let $x$ be
the largest element of $S_0$ with $x\le i_1$ (such elements exist, as
$a<i_1$). Then $x+i<i_1+b\le 2b\le k$, so $C_0(i)$ gives $x+i\in S_0$. By the
maximality of $x$, $x+i>i_1$, so $x+i$ lies in the interval
\eqref{eq:nao-evengap}, a contradiction.

This proves the claim. The same claim for the even element $k-\sigma$, with
\autoref{lem:nao-parity}(b) and \autoref{lem:nao-mirror}, gives a surviving
tap reading $0$ at every index for $Q_\sigma$. With
\autoref{lem:nao-parity}(a), $(\sigma,Q_\sigma)$ is a $\sigma$-bump.

For the last sentence: either $(\sigma,Q_\sigma)$ is a $\sigma$-bump for
some odd $\sigma$, or it fails at $\sigma_1=b$ and then $(a,Q_a)$ is one.
\end{proof}

\subsection{Odd minimum: where the failures sit}

From now on $a=\min S$ is odd and $S_0\ne\emptyset$, so that $b=\min S_0$
in the notation above. Then $a<b\le k-b$, $\max S_1=k-a$ and
$\max S_0=k-b$. Note that $a+b<k$, and that an odd number at most $b$
is less than $b$.

\begin{lemma}[localisation]
\label{lem:nao-loc}
\begin{enumerate}
\item[(a)] If $\sigma\in S$ is even and $i$ is silent for $Q_\sigma$, then
$C_0(i)$ holds and either $i=a$, or $i$ is odd with $a<i<b$.
\item[(b)] If $\sigma\in S$ is odd and $i$ is silent for $Q_\sigma$, then
$C_1(i)$ holds and either $i$ is odd with $i\le a$, or $i$ is even with
$a<i\le b$.
\end{enumerate}
\end{lemma}

\begin{proof}
The containments are \autoref{lem:nao-parity}(e). (a) An even $i<a$ is
excluded by \autoref{lem:nao-ends}(ii), as $i<a<b=\min S_0$; an odd $i<a$
by \autoref{lem:nao-ends}(i), as $i\equiv a$. If $i>a$, then
\autoref{lem:nao-dichotomy} makes $i$ odd with $i\le b$, hence $i<b$.
(b) An even $i<a$ is excluded by \autoref{lem:nao-ends}(i), as
$i\equiv a+1$; and $i=a$ is odd. If $i>a$, then \autoref{lem:nao-dichotomy}
makes $i$ even with $i\le b$.
\end{proof}

So a silent index of the even class sits at the end $i=a$ or in the odd part
of the inner range $(a,b)$, and a silent index of the odd class at an odd
end $i\le a$ or in the even part of the inner range $(a,b]$. When both classes fail,
three of the four pairings are impossible, and the fourth is rigid.

\begin{lemma}[no inner failure of the even class]
\label{lem:nao-noband}
Let $i_0$ be odd with $a<i_0<b$ and $C_0(i_0)$, and let $i_1$ satisfy
$C_1(i_1)$, with either $i_1$ odd and $i_1\le a$, or $i_1$ even and
$a<i_1\le b$. This is impossible.
\end{lemma}

\begin{proof}
By \autoref{lem:nao-gap} with $c=0$, $S$ has no odd element in
$(i_0,\ i_0+b)$.

If $i_1$ is odd and $i_1\le a$, then $b+i_1<2b\le k$ is odd, so $C_1(i_1)$
puts it in $S$; but $i_0<b<b+i_1<b+i_0$.

If $i_1$ is even, let $x$ be the largest odd element of $S$ with $x\le i_0$
(such elements exist, as $a<i_0$). Then $x+i_1<2b\le k$ is odd, so $x+i_1\in
S$ by $C_1(i_1)$, and $x+i_1>i_0$ by the maximality of $x$. It cannot lie in
the gap, so $x+i_1\ge i_0+b$, which with $x\le i_0$ forces $i_1=b$. But
then $a+b\in S$ in the same way, and $i_0<a+b<i_0+b$.
\end{proof}

The pairing of the two end failures is excluded by primitivity, and this is
where the full support appears as the one obstruction. This proof and the
next use an elementary observation. Say that a finite set $W$ with largest
element $w$ is \emph{closed under the step} $r>0$ if $y+r\in W$ for every
$y\in W$ with $y\ne w$. Then, for every $y\in W$, the sequence
$y,y+r,y+2r,\dots$ stays in $W$ until it reaches $w$, which it must do
because $W$ is bounded by $w$; so $r\mid w-y$, and $y,y+r,\dots,w$ all lie in
$W$.

\begin{lemma}[two end failures]
\label{lem:nao-endend}
Let $S$ be primitive, and suppose $C_0(a)$ and $C_1(i_1)$ hold with $i_1$
odd and $i_1\le a$. Then $S$ is the full support.
\end{lemma}

\begin{proof}
Put $r=a+i_1$. Composing the two containments makes each class closed under
a step, and the two steps turn out to be the same.

\emph{The odd class has step $r$.} Let $y\in S_1$ with $y<k-a$. Then $a+y<k$
is even, so $a+y\in S_0$ by $C_0(a)$, whence $a+y\le k-b$. Then
$a+y+i_1<k$ (as $i_1<b$) is odd, so $y+r\in S_1$ by $C_1(i_1)$. So $S_1$
is closed under the step $r$, with top $k-a$; every $y\in S_1$ has
$r\mid k-a-y$, and in particular $r\mid k-2a$.

\emph{$b=r$.} The position $k-b+i_1<k$ is odd, so it lies in $S_1$ by
$C_1(i_1)$, and $k-b+i_1\le k-a$. If the inequality were strict, $C_0(a)$
would put the even position $a+i_1+k-b<k$ in $S_0$, above
$\max S_0=k-b$. So $k-b+i_1=k-a$, that is $b=r$.

\emph{The even class has step $b$.} Let $y\in S_0$ with $y<k-b$. Then
$y+i_1<k$ is odd, so $y+i_1\in S_1$ by $C_1(i_1)$; and
$y+i_1<k-b+i_1=k-a$, so $a+y+i_1<k$ is even and lies in $S_0$ by $C_0(a)$.
That is, $y+b\in S_0$. So $S_0$ is closed under the step $b$, with top
$k-b$, and every $y\in S_0$ has $b\mid k-b-y$; at $y=b$ this gives
$b\mid k$.

\emph{Collapse.} Now $b\mid k$ and $b=r\mid k-2a$, so $b\mid 2a$; since
$b>a$, $b=2a$, and $a\mid k$. Every odd $s\in S$ has $2a\mid k-a-s$ and
every even $s\in S$ has $2a\mid k-2a-s$, so $a$ divides every element of
$S\cup\{k\}$, and primitivity gives $a=1$. Then $r=2$, and the chain from $a=1$
puts every odd number in $[1,k-1]$ in $S$; and for even $y\in(0,k)$,
$y-1\in S$ is odd with $1+(y-1)<k$, so $y\in S$ by $C_0(1)$. So $S$ is full.
\end{proof}

\begin{lemma}[end against inner range]
\label{lem:nao-endband}
Suppose $C_0(a)$ and $C_1(i_1)$ hold with $i_1$ even and $a<i_1\le b$.
Then $i_1=b$, $b\mid k-2a$, and
$S_1=\{a,\ a+b,\ a+2b,\ \dots,\ k-a\}$.
\end{lemma}

\begin{proof}
By \autoref{lem:nao-gap} with $i=a$ and $c=0$, $S$ has no odd element in
$(a,a+b)$. The position $a+i_1\le a+b<k$ is odd, so it lies in $S$ by
$C_1(i_1)$, hence $a+i_1\ge a+b$ and $i_1=b$. Now let $y\in S_1$ with
$y<k-a$. Then $k-y\in S_1$ exceeds $a$, so $k-y\ge a+b$, that is
$y+b\le k-a$; and $y+b$ is odd, so it lies in $S_1$ by $C_1(b)$.
So $S_1$ is closed under the step $b$, with top $k-a$: every element of
$S_1$ is $\equiv k-a\pmod b$, and the whole progression from $a$ to $k-a$
lies in $S_1$. Hence $b\mid k-2a$ and $S_1$ is that progression.
\end{proof}

\begin{theorem}[the both-failing structure]
\label{thm:nao-struct}
Let $k$ be even and $S$ primitive and palindromic, not the full support,
with $a=\min S$ odd and $S_0\ne\emptyset$; put $b=\min S_0$. Suppose that
$(\sigma,Q_\sigma)$ is not a $\sigma$-bump for any $\sigma\in S$. Then
\[
k=2a+Mb\quad\text{for an integer } M\ge1,\qquad a<b<2a,
\]
\[
S_1=\{a+mb:\ 0\le m\le M\},\qquad
2a+mb\in S\ \text{ for } 0\le m<M .
\]
\end{theorem}

\begin{proof}
By \autoref{lem:nao-parity}(d) applied at $\sigma=b$, some index $i_0$ is
silent for $Q_{\sigma_0}$ with $\sigma_0$ even, and at $\sigma=a$ some
index $i_1$ is silent for $Q_{\sigma_1}$ with $\sigma_1$ odd.
\autoref{lem:nao-loc} locates both, and \autoref{lem:nao-noband} rules out
$i_0\ne a$. So $C_0(a)$ holds.

Next, let $\sigma\in S$ be odd and let $i$ be any index silent for
$Q_\sigma$. By \autoref{lem:nao-loc}(b) and \autoref{lem:nao-endend} (with
$C_0(a)$ and since $S$ is not full), $i$ is even with $a<i\le b$, and then
\autoref{lem:nao-endband} gives
\begin{equation}
i=b,\qquad b\mid k-2a,\qquad S_1=\{a,a+b,\dots,k-a\}.
\label{eq:nao-band}
\end{equation}
Applying this to $i_1$ gives the description of $S_1$, with $M=(k-2a)/b$.
Here $M\ge1$ because $2a<k$ by \eqref{eq:nao-ab}. For $0\le m<M$, the odd
element $a+mb$ has $a+(a+mb)\le 2a+(M-1)b<k$, so $C_0(a)$ puts
$2a+mb$ in $S$. At $m=0$, $2a\in S_0$, so $b\le 2a$.

It remains to exclude $b=2a$; suppose it holds. If $M=1$ then $k=4a$,
$S_1\subseteq\{a,3a\}$ and $S_0\subseteq[b,k-b]=\{2a\}$, so $a$ divides
every element of $S\cup\{k\}$. Primitivity gives $a=1$, $k=4$, and
$S\supseteq\{a,b,a+b\}=\{1,2,3\}$ is full, against the hypothesis. If
$M\ge2$, then $3a=a+b\in S_1$, and by \autoref{lem:nao-parity}(d) some
index $i_3$ is silent for $Q_{\sigma_3}$ with $\sigma_3\in\{3a,k-3a\}$. By
\eqref{eq:nao-band} $i_3=b=2a$. But there the tap $k-a$ is wrapped and reads the
pin $k+a$, of residue $a$, and $a\ne\sigma_3$ (as $k-3a=a$ would mean $k=4a$,
that is $M=1$). So the tap survives and reads $1$ by
\eqref{eq:read-pins}, and $i_3$ is not silent.
\end{proof}

\subsection{Carriers for the both-failing supports}

The structure of \autoref{thm:nao-struct} leaves the even elements of $S$
strictly inside $(b,k-b)$ largely undetermined. Both carriers below are
built to be insensitive to them.

\begin{lemma}[the case $M=1$]
\label{lem:nao-bfone}
Let $a$ be odd and $b$ even with $a<b<2a$, put $k=2a+b$, and let
$S\subseteq\{1,\dots,k-1\}$ be palindromic with $a,b\in S$, with every odd
element in $\{a,a+b\}$ and every even element in $[b,2a]$. Then $\sigma=a$
with
\[
Q=\{a,\ 3a,\ b-a\}\cup\{\,y\in\{1,\dots,k-1\}:\ y \text{ even},\ y\notin S\,\}
\]
is a $\sigma$-bump.
\end{lemma}

\begin{proof}
By palindromy $a+b=k-a$ and $2a=k-b$ lie in $S$. The positions $3a<k$ and
$b-a>0$ are odd and differ from $a$ and $a+b$ (since $b\ne2a$), so
$Q\cap S=\{a\}$, which is (B1). An odd position lies in $Q$ exactly when it
is one of $a$, $3a$, $b-a$; an even position lies in $Q$ exactly when it is
free, and in particular every even position in $(0,b)\cup(2a,k)$ lies in $Q$.

For (B2) it suffices, at each index $i$, to name two surviving taps among
$a$, $b$, $2a$, $a+b$ whose reads differ. The indices fall into twelve
ranges, six of even and six of odd indices, cut out by the numbers $a$, $b$,
$2a$, $a+b$, $2b$ and $3a$, and in each range the same two taps serve. The
table of these taps, and the check that they survive and read different
values, are in \autoref{app:checks-bfone}.
\end{proof}

For $M\ge2$ the carrier is the \emph{band carrier}: $a$ together with every
free position strictly between $a$ and $k-a$.

\begin{lemma}[the case $M\ge2$]
\label{lem:nao-bfshadow}
Let $S$ have the structure of \autoref{thm:nao-struct} with $M\ge2$, and let
\[
\mathcal B=\{\,y\notin S:\ a<y<k-a\,\}.
\]
Then $\sigma=a$ with $Q=\{a\}\cup\mathcal B$ is a $\sigma$-bump.
\end{lemma}

The idea is that at an index $i>a$ the tap $k-a$ reads the position $i-a$,
wrapped, and so reads $1$ unless $i-a$ lies in $\mathcal B$; when it does, the
gaps of $S$, which are shorter than $a$, supply an element $j$ of $S$ just
below $i$, and the tap $k-j$ reads a position below $\mathcal B$. The mirror then
gives the reads of $0$.

\begin{proof}
\begin{itemize}
\item[(N1)] For $1\le m\le M$, $mb=k-(2a+(M-m)b)\in S$, by palindromy and
the last clause of the structure.
\item[(N2)] An odd $y\in[a,k-a]$ with $y-a$ not a multiple of $b$ is free,
since $S_1=\{a+mb\}$.
\item[(N3)] If $a<x\le k-a$, some $j\in S$ has $x\le j<x+a$.
\end{itemize}
For (N3), I first show that every $s\in S$ with $s<k-a=a+Mb$ has an
element of $S$ in $(s,s+a]$. Write $s=qb+r$ with $0\le r<b$. If $r<a$,
then $q\le M$ (as $qb\le s<a+Mb<(M+1)b$), and $a+qb\in S$ lies in
$(s,s+a]$. If $r\ge a$, then $qb+a\le s<a+Mb$ gives $q<M$, and
$(q+1)b\in S$ by (N1); it lies in $(s,s+a]$ because
$(q+1)b-s=b-r\le b-a<a$. Now, given $x$, let $s$ be the largest element
of $S$ below $x$ (there is one, $a$). Then $s<k-a$, so some $j\in S$ lies in
$(s,s+a]$, and $j\ge x$ by the maximality of $s$, while $j\le s+a<x+a$.

The set $\mathcal B$ consists of free positions, so $Q\cap S=\{a\}$, which is (B1).
By palindromy $\mathcal B$ is symmetric under $y\mapsto k-y$, so $Q'=\{k-a\}\cup \mathcal B$ is
the reflection of $Q$.

\emph{Reads of $1$.} Let $z\in\{a,k-a\}$ and $Q^z=\{z\}\cup \mathcal B$; I show that
every index $i$ has a surviving tap reading $1$ for $(z,Q^z)$. As
$z\in S$, $z\notin \mathcal B$, and a tap whose position lies in $\mathcal B$ survives.
\begin{itemize}
\item If $i<a$ and $i$ is even, the tap $a$ reads $i+a$, which is odd, lies in
$(a,2a)\subset(a,k-a)$, and has $(i+a)-a=i\in(0,b)$. By (N2) it lies in
$\mathcal B$, and the tap, being unwrapped, reads $1$.
\item If $i<a$ and $i$ is odd, the tap $b$ reads $i+b$, which is odd, lies
in $(b,a+b)\subset(a,k-a)$, and has $(i+b)-a\in(b-a,b)\subset(0,b)$.
So it lies in $\mathcal B$ and reads $1$.
\item If $i=a$, the tap $2a\in S$ reads $3a$, which is odd, has
$3a-a=2a\in(b,2b)$, and lies in $(a,k-a)$, since $2a<2b\le Mb$. So it
lies in $\mathcal B$ and reads $1$.
\item If $i>a$, the tap $k-a$ is wrapped, with residue $y=i-a\in(0,k-a)$. If
$y\notin \mathcal B$ and $y\ne z$, it survives and reads $1$.
\item If $i>a$ and $i-a=z$, then $z\ne k-a$, so $z=a$ and $i=2a$. The tap $k-b$ is
wrapped (as $2a>b$), with residue $2a-b\in(0,a)$, which is neither in $\mathcal B$
nor $z$; so it survives and reads $1$.
\item If $i>a$ and $i-a\in \mathcal B$, then $a<i-a<k-a$, and (N3) with $x=i-a$ gives $j\in S$
with $i-a\le j<i$. The tap $k-j$ is wrapped, with residue $y=i-j\in(0,a]$,
which is not in $\mathcal B$. If $y=z$, then $z\le a<k-a$ forces $z=a$, so
$j=i-a\in \mathcal B$, which is impossible as $j\in S$. So $y\ne z$, and the tap
survives and reads $1$.
\end{itemize}
With $z=a$ this gives a read of $1$ at every index for $(a,Q)$. With $z=k-a$
at the index $k-i$, \autoref{lem:nao-mirror} gives a surviving read of $0$
at $i$ for $(a,Q)$. So (B2) holds.
\end{proof}

\begin{corollary}[the both-failing supports]
\label{cor:nao-bothfailing}
Under the hypotheses of \autoref{thm:nao-struct}, $S$ carries a $\sigma$-bump
at $\sigma=a$.
\end{corollary}

\begin{proof}
If $M\ge2$ this is \autoref{lem:nao-bfshadow}. If $M=1$, then $k=2a+b$, the
odd elements of $S$ are $a$ and $a+b$, and every even element $s$ has
$b\le s$ and $b\le k-s$, so $s\in[b,2a]$. This is
\autoref{lem:nao-bfone}.
\end{proof}

\subsection{Proof of \autoref{thm:notallodd}}

\begin{proof}[Proof of \autoref{thm:notallodd}]
Let $a=\min S$. If $a$ is even, \autoref{thm:nao-evenmin} gives a
$\sigma$-bump on a parity carrier. If $a$ is odd, either
$(\sigma,Q_\sigma)$ is a $\sigma$-bump for some $\sigma\in S$, or $S$
satisfies the hypotheses of \autoref{thm:nao-struct}, since it is primitive,
palindromic, not full and has an even element; and then
\autoref{cor:nao-bothfailing} gives a $\sigma$-bump at $\sigma=a$.
\end{proof}

By \autoref{thm:sigmabump}, no equation \eqref{eq:familya} whose support
satisfies the hypotheses of \autoref{thm:notallodd} is globally periodic,
for either value of $\theta$.
